\documentclass[11pt]{amsart}
\usepackage{amsmath,amssymb,amsthm,booktabs}
\usepackage[margin=1.05in]{geometry}
\newtheorem{theorem}{Theorem}[section]
\newtheorem{proposition}[theorem]{Proposition}
\newtheorem{lemma}[theorem]{Lemma}
\newtheorem{corollary}[theorem]{Corollary}
\newtheorem{remark}[theorem]{Remark}
\newcommand{\eu}{\mathrm{e}}
\newcommand{\D}{\mathcal D}
\newcommand{\Def}{\mathrm{Def}}
\title[Analytic Zone C for $m\in\{9,11,13\}$]{An analytic Zone-C bound for the
short cycles $m\in\{9,11,13\}$\\ (the $\xi\le1$ part of Region~II, all $e$)}
\author{Claude (v2, 2026-07-17)}
\begin{document}
\maketitle

\section{Statement and specialization}\label{zc9:sec:setup}

This note proves, for each fixed odd $m\in\{9,11,13\}$, the scalar Region-II
inequality on the whole $\xi\le1$ (Zone-C) slice of the domain.  It specializes
the verified analytic Zone-C proof for $m\ge15$
(\texttt{proof\_zoneC\_analytic\_v2.tex}, hereafter \textbf{[ZC]}) to the fixed
exponent $n=m-2\in\{7,9,11\}$.  The single new analytic ingredient is a
\emph{mean-value Case-II row bound} (\S\ref{zc9:sec:mvt}) that removes the
$m$-uniform supremum loss of [ZC] for the small exponents; with it, every range
closes on a \textbf{fixed, fully displayed} interval cover.  There is no
adaptive or greedy subdivision anywhere in this note.  Every displayed constant
is an exact rational rounded in the stated safe direction and is re-checked by
the companion script \texttt{validate\_zoneC\_91113\_v2.py} (dense float grids
over the exact domains, \texttt{Fraction} arithmetic at the endpoints; exit
$0$).

We use the $(e,\kappa)$ chart of the Region-II manuscript
(\texttt{paper\_region2\_v2.tex}, hereafter \textbf{[R2]}) exactly as in [ZC]:
\[
 \alpha=\tfrac{1-e}2,\ \ d=\alpha-q,\ \ \delta=\tfrac{1-3e}6,\ \
 \kappa=\tfrac de,\ \ p=1-q,\ \ L=\sqrt{pq-\alpha^2},\ \
 x=\tfrac\alpha p,\ s=\tfrac qp,\ y=\tfrac Lp,\ \ell=\tfrac L\alpha,
\]
\[
 f=\alpha-L,\qquad \xi=\tfrac{4\alpha^2d}{e^2},\qquad
 R_m=\alpha^m+L^m-pq^{m-1},\qquad n:=m-2 .
\]
The target is the Huber inequality of Lemma~\emph{huber} of [R2]:
\begin{equation}\label{zc9:eq:target}
 R_m\ \le\ C_m\,\psi(\xi,\rho),\qquad
 C_m=\tfrac{\sqrt{2\alpha}}{4\alpha^2}B_mf e^2,\quad
 \psi=\max_{0\le\lambda\le1}\Bigl(\lambda\xi-\tfrac{\lambda^2}{4(\rho+\lambda)}\Bigr),
\end{equation}
with $A_m=2L^{n}+m\,k_m(\alpha)$, $B_m=2L^{n}+m\,k_m(L)$,
$k_m(\mu)=\tfrac{p^{m-1}-\mu^{m-1}}{p+\mu}$ and
$\rho=\tfrac{A_m}{B_m}\tfrac{\sqrt\alpha}{2\sqrt2 f}$.  Write
$\Def_m=R_m/(\alpha^3p^{\,n})$.

\begin{theorem}\label{zc9:thm:main}
Let $m\in\{9,11,13\}$ and let $(q,\alpha)$ be any Region-II point with $\xi\le1$
(equivalently $d\le e^2/(1-e)^2$) and $q>\tfrac13$.  Then
\eqref{zc9:eq:target} holds.
\end{theorem}

The proof partitions $e\in(0,\tfrac13)$ into the pieces of
Table~\ref{zc9:tab:split}; the strip end $e^\star_m$ equals the corner start for
$m\in\{11,13\}$, while for $m=9$ a short \emph{sliver} bridges the two.
\begin{table}[h]\centering\small
\caption{Zone-C partition.  $e^\star_m=\tfrac13-2\Delta_m$ with
$\Delta_{13}=\tfrac{13}{600},\Delta_{11}=\tfrac9{600},\Delta_9=\tfrac7{600}$.}
\label{zc9:tab:split}
\begin{tabular}{llll}
\toprule
piece & range & method & where\\
\midrule
Zone A & $0<e\le\tfrac1{60}$ & $R_m\le0$ & Cor.~\ref{zc9:cor:za}\\
strip  & $\tfrac1{60}\le e\le \bar e_m$ & fixed $\varphi$-row table & Prop.~\ref{zc9:prop:strip}\\
sliver & ($m=9$) $\tfrac{29}{100}\le e\le\tfrac{31}{100}$ & fixed $(e,t)$ cell table & Prop.~\ref{zc9:prop:sliver}\\
corner & $e^\star_m\le e<\tfrac13$ & dual-certificate chain & Prop.~\ref{zc9:prop:wide}\\
\bottomrule
\end{tabular}
\end{table}
Here $\bar e_{13}=e^\star_{13}=\tfrac{29}{100}$,
$\bar e_{11}=e^\star_{11}=\tfrac{91}{300}$, and
$\bar e_{9}=\tfrac{29}{100}$ with $e^\star_9=\tfrac{31}{100}$.

\subsection*{What is inherited verbatim from [ZC]}
[ZC] states (its closing remark) that it uses $m\ge15$ \emph{only} through
$n\ge13$ (the exponent in $G_2$), $1-y^{m-1}\ge1-y^{14}$, and the constant $15$
in its Lemma~\emph{master}.  Replacing $13\rightsquigarrow n=m-2$,
$14\rightsquigarrow m-1$ and $15\rightsquigarrow m$ throughout leaves the
following results and their proofs intact for every fixed $m\ge9$ ($n\ge7$); we
quote them by the [ZC] names, with
\[
 \lambda_x=\ln\tfrac p\alpha,\quad \Delta=\ln\tfrac\alpha q,\quad
 \lambda_s=\lambda_x+\Delta,\quad G_2=\ell^2\bigl(1-(y/s)^{n}\bigr),\quad
 G=\ln(1+G_2),\quad n_g=\tfrac G\Delta .
\]
\begin{itemize}
\item \textbf{[ZC] Lemma~\emph{majorant}} (three-geometric identity of [R2]
 Lemma~\emph{threegeo}(i) + $s^n-y^n\ge s^nG_2$):
 $\alpha\,\Def_m=x^n-s^n-\ell^2(s^n-y^n)\le\varphi(n):=x^n-(1+G_2)s^n$.
\item \textbf{[ZC] Lemma~\emph{sup}}: $\varphi(t)=x^t-(1+G_2)s^t$ is unimodal with
 $\sup_t\varphi=(\Delta/\lambda_s)h\,\eu^{-w}$, $w:=\lambda_xG/\Delta$,
 $h:=\exp(-\tfrac{\lambda_x}\Delta\ln(1+\tfrac\Delta{\lambda_x}))\in(\eu^{-1},1)$,
 $h\le\eu^{-1+v/2}$, $v=\Delta/\lambda_x$.
\item \textbf{[ZC] Lemma~\emph{scalar}}, \textbf{Lemma~\emph{payfloor}} and
 \textbf{Lemma~\emph{master}}: with
 $c_e=\tfrac{2(1-x^{m-1})(1-y^{m-1})}{(1+x)(1+y)}$,
 $K_2=\tfrac{\sqrt{2\alpha}(1-y^{m-1})}{2\alpha^3(1+y)}$, the payment floors are
 $\Def_m^{-1}C_m\psi/(\alpha^{-3}p^{-n})\ge m c_e\kappa^2$ if $2\rho\xi\le1$ and
 $\ge mK_2fd$ if $2\rho\xi>1$; and, if $R_m>0$, \eqref{zc9:eq:target} follows
 from the applicable one of the \emph{Case-I} master inequalities
 \begin{align}
  &\textup{(deep,}\ 2\rho\xi\le1,\ n_g\ge m)&
   h\eu^{-w}e^2&\le\lambda_x\alpha c_e q^2G,\label{zc9:eq:SCdeep}\\
  &\textup{(shallow,}\ 2\rho\xi\le1,\ n_g<m)&
   h\eu^{-w}e^2&\le m\lambda_x\alpha c_e q^2\Delta.\label{zc9:eq:SCsh}
 \end{align}
\item \textbf{[ZC] Lemma~\emph{bracket}}: over an interval $E=[e_L,e_R]$ with
 $d\le d_{\max}(e)=\min(e^2/(1-e)^2,\delta)$, the monotone-endpoint constants
 $\bar q,D,\Lambda,\lambda_L,\widehat G_2,Y,\gamma,a,w^-,v^+,h^+,x_U,y_U,c^-,
 f^-,K^-$ of [ZC] Lemma~\emph{bracket} (recorded in \S\ref{zc9:sec:strip})
 satisfy $q\ge\bar q$, $\Delta\le\Lambda$, $\lambda_x\ge\lambda_L$,
 $\lambda_xG\ge a$, $w\ge w^-$, $h\le h^+$, $c_e\ge c^-$, $f\ge f^-$,
 $K_2\ge K^-$, $x\le x_U$, $y\le y_U$.
\end{itemize}
The oddness of $m$ and $\alpha\le r(q)$ enter only through $L^2>0$,
$0\le\ell<1$, $0<y<s$, $L<q$ (i.e.\ [R2] Lemma~\emph{frontier-gap}, valid on all
of $\D$).  This note replaces only the \emph{Case-II} treatment of [ZC] (its
master inequality \eqref{zc9:eq:SCII-old} below and the resulting row bound
$\Psi_{\mathrm{II}}$): for $m\le13$ the $m$-uniform supremum
$\varphi(n)\le\sup_t\varphi$ overshoots by a factor $\approx4$ on the Case-II
ridge, and we recover it by the mean-value bound of \S\ref{zc9:sec:mvt}.  The
[ZC] Case-II master inequality, retained as one of two admissible Case-II
sufficient conditions, is
\begin{equation}\label{zc9:eq:SCII-old}
 \textup{(Case II, }2\rho\xi>1):\qquad
 h\eu^{-w}\ \le\ m\lambda_x\alpha K_2 f q .
\end{equation}

\section{Zone A: $e\le\tfrac1{60}$}\label{zc9:sec:za}

\begin{corollary}\label{zc9:cor:za}
For every $m\in\{9,11,13\}$ and every Region-II point with $\xi\le1$ and
$e\le\tfrac1{60}$ one has $R_m\le0$; hence \eqref{zc9:eq:target} holds trivially
($C_m,\psi\ge0$).
\end{corollary}

\begin{proof}
This is [R2] Cor.~\emph{gate-smalle} at the fixed $m$: [R2]
Prop.~\emph{smalle} (valid for all odd $3\le m\le13$) gives $R_m\le0$ on
$\{e\le\tfrac1{60},\ \xi\le1\}$.  Part~7 of the validation script confirms
$\max R_m<0$ there (grid maxima $-3.9\cdot10^{-7},-9.7\cdot10^{-8},
-2.4\cdot10^{-8}$ for $m=9,11,13$).
\end{proof}

Because of Corollary~\ref{zc9:cor:za} the strip may start at $e=\tfrac1{60}$.

\section{A mean-value Case-II row bound}\label{zc9:sec:mvt}

The only structural change from [ZC] is to replace, in Case~II, the majorant
$\alpha\Def_m\le\varphi(n)\le\sup_t\varphi$ (which loses a factor
$\approx\!4$ at $n=7$) by the sharper mean-value majorant already used by [ZC]
in its wide corner.

\begin{lemma}[mean-value defect]\label{zc9:lem:mvt}
On all of $\D$, $\alpha\,\Def_m\le\dfrac{n\,x^{\,n-1}}p\,N$ where
$N:=d-\ell^{\,m-1}(q-L)$, and moreover $0<N\le d$ whenever $R_m>0$.
\end{lemma}

\begin{proof}
By [ZC] Lemma~\emph{majorant}, $\alpha\Def_m=(x^n-s^n)-\ell^2(s^n-y^n)$.  The
mean value theorem gives $c'\in(s,x)$, $c\in(y,s)$ with $x^n-s^n=n c'^{\,n-1}
(x-s)\le n x^{n-1}(x-s)$ and $s^n-y^n=n c^{\,n-1}(s-y)\ge n y^{n-1}(s-y)$.
Since $\ell^2>0$,
\[
 \alpha\Def_m\ \le\ n x^{n-1}(x-s)-\ell^2 n y^{n-1}(s-y).
\]
Now $x-s=(\alpha-q)/p=d/p$, $s-y=(q-L)/p$, and $y=\ell x$ gives
$\ell^2 y^{n-1}=\ell^{\,n+1}x^{n-1}=\ell^{\,m-1}x^{n-1}$ (as $m-1=n+1$).  Hence
$\alpha\Def_m\le\tfrac{n x^{n-1}}p\bigl(d-\ell^{\,m-1}(q-L)\bigr)
=\tfrac{n x^{n-1}}p N$.  Finally $q>L$ on $\D$ ($L<q\iff q>\tfrac13$) and
$\ell>0$, so $\ell^{\,m-1}(q-L)\ge0$ and $N\le d$; and $R_m>0$ forces
$\alpha\Def_m>0$, whence $N>0$.  (Part~4 of the script verifies both
inequalities on a dense grid.)
\end{proof}

The point is that $N\le d$ makes the Case-II comparison \emph{$d$-free}: the
defect and the floor $mK_2fd$ are both proportional to $d$.

\begin{lemma}[Case-II master, MVT form]\label{zc9:lem:mvtII}
Let $R_m>0$ and $2\rho\xi>1$.  Then \eqref{zc9:eq:target} holds as soon as
\begin{equation}\label{zc9:eq:SCII-mvt}
 \frac{n\,x^{\,n-1}}{p\,\alpha}\ \le\ m\,K_2\,f .
\end{equation}
\end{lemma}

\begin{proof}
By [ZC] Lemma~\emph{payfloor}, $2\rho\xi>1$ gives
$C_m\psi/(\alpha^3p^{n})\ge mK_2fd$, i.e.\ it suffices that $\Def_m\le mK_2fd$.
By Lemma~\ref{zc9:lem:mvt}, $\Def_m\le\tfrac{n x^{n-1}}{p\alpha}N\le
\tfrac{n x^{n-1}}{p\alpha}d$.  Dividing by $d>0$ turns $\Def_m\le mK_2fd$ into
\eqref{zc9:eq:SCII-mvt}.
\end{proof}

\begin{lemma}[assembled Case-II row bound]\label{zc9:lem:mvtrow}
Over an interval $E=[e_L,e_R]$ with $d\le d_{\max}$, put
$p^-:=\tfrac{1+e_L}2$ (so $p\ge p^-$) and take $x_U,f^-,K^-$ from [ZC]
Lemma~\emph{bracket}.  If
\[
 \Psi_{\mathrm{II}}^{\mathrm{mvt}}(E):=
 \frac{n\,x_U^{\,n-1}}{p^-\,\alpha(e_R)\,m\,K^-\,f^-}\ \le\ 1,
\]
then \eqref{zc9:eq:SCII-mvt} holds throughout $E\cap\{2\rho\xi>1,\ R_m>0\}$.
\end{lemma}

\begin{proof}
$x\le x_U$ and $n\ge1$ give $n x^{n-1}\le n x_U^{\,n-1}$.  Also $p\ge1-\alpha=
\tfrac{1+e}2\ge p^-$, $\alpha\ge\alpha(e_R)$, $f\ge f^-$, and $K_2\ge K^-$
(the last three are [ZC] Lemma~\emph{bracket}(a),(e), unchanged).  Hence
$\tfrac{n x^{n-1}}{p\alpha}\le\tfrac{n x_U^{\,n-1}}{p^-\alpha(e_R)}$ and
$mK_2f\ge m K^-f^-$, so
$\Psi_{\mathrm{II}}^{\mathrm{mvt}}(E)\le1$ implies \eqref{zc9:eq:SCII-mvt}.
\end{proof}

Since \eqref{zc9:eq:SCII-old} and \eqref{zc9:eq:SCII-mvt} are \emph{both}
sufficient for the Case-II target, the effective Case-II row bound on $E$ is
\begin{equation}\label{zc9:eq:PsiII}
 \Psi_{\mathrm{II}}(E):=\min\bigl(\Psi_{\mathrm{II}}^{\mathrm{[ZC]}}(E),\
 \Psi_{\mathrm{II}}^{\mathrm{mvt}}(E)\bigr),
\end{equation}
where $\Psi_{\mathrm{II}}^{\mathrm{[ZC]}}(E)=h^+\eu^{-a/\Lambda}/
(m\lambda_L\alpha(e_R)K^-f^-\bar q)$ is [ZC]'s row bound for
\eqref{zc9:eq:SCII-old}.  The [ZC] form is tighter at small $e$ (where $N\ll d$,
so the $d$-free bound is loose); the MVT form is tighter on the ridge.
No $d$-lower-bound (the ``$d_{\mathrm{lo}}$ upgrade'' of the previous draft) is
needed.

\section{The strip $\tfrac1{60}\le e\le\bar e_m$}\label{zc9:sec:strip}

Fix $E=[e_L,e_R]$ with $d\le d_{\max}(e)$.  The [ZC] Lemma~\emph{bracket}
constants (elementary evaluations at the rational endpoints;
$\alpha(u)=\tfrac{1-u}2$, $\delta(u)=\tfrac{1-3u}6$) are
\[
 \bar q=\max\bigl(\tfrac13,\alpha(e_R)-\tfrac{e_R^2}{(1-e_R)^2}\bigr),\ \
 D=\min\bigl(\tfrac{e_R^2}{(1-e_R)^2},\delta(e_L)\bigr),\ \
 \Lambda=\tfrac D{\bar q},\ \
 \lambda_L=\ln\tfrac{1+e_L}{1-e_L},\ \
 \widehat G_2=\tfrac{2e_R}{1-e_R},
\]
\[
 Y=\tfrac{\sqrt{e_R(1-e_R)/2}}{\bar q},\quad
 \gamma=\Bigl[\tfrac2{1-e_L}-\tfrac{D(e_R+D)}{\alpha(e_R)^2e_L}\Bigr]
 (1-Y^n)\tfrac{\ln(1+\widehat G_2)}{\widehat G_2},\quad
 a=\lambda_L\gamma e_L,
\]
\[
 w^-=2\gamma\max\bigl((1-e_R)^2\bar q,\tfrac{e_L^2}{3\delta(e_L)}\bigr),\quad
 v^+=\min\bigl(\tfrac{e_R}{2(1-e_R)^2\bar q},\tfrac{3\delta(e_L)}{\lambda_L}\bigr),\quad
 h^+=\eu^{-1+v^+/2},
\]
\[
 x_U=\tfrac{1-e_L}{1+e_L},\ \
 y_U=\tfrac{\sqrt{e_R(1-e_R)/2}}{\alpha(e_R)+e_R},\ \
 c^-=\tfrac{2(1-x_U^{m-1})(1-y_U^{m-1})}{(1+x_U)(1+y_U)},\ \
 f^-=\alpha(e_R)-\sqrt{\tfrac{e_R(1-e_R)}2},
\]
\[
 K^-=\tfrac{\sqrt{2\alpha(e_R)}(1-y_U^{m-1})}{2\alpha(e_L)^3(1+y_U)},\qquad
 p^-=\tfrac{1+e_L}2 .
\]
As in [ZC] Lemma~\emph{rows}, the deep and shallow (Case-I) sufficient
conditions \eqref{zc9:eq:SCdeep}--\eqref{zc9:eq:SCsh} reduce, via
Lemma~\emph{bracket}, to $\Psi_{\mathrm{deep}}(E)\le1$ and (when the shallow case
is nonempty, i.e.\ $\Lambda>\gamma e_L/m$) $\Psi_{\mathrm{sh}}(E)\le1$, where
\[
 \Psi_{\mathrm{deep}}=\frac{h^+\eu^{-\max(w^-,m\lambda_L)}}
 {2\gamma\,\alpha(e_R)\,c^-\,\bar q^{\,2}},\qquad
 \Psi_{\mathrm{sh}}=\frac{h^+(e_R/2)\eu^{-a/\Delta^*}}
 {m\,\alpha(e_R)\,c^-\,\bar q^{\,2}\,\Delta^*},\quad
 \Delta^*=\min\bigl(\Lambda,\max(a,\tfrac{\gamma e_L}m)\bigr),
\]
and the Case-II condition reduces to $\Psi_{\mathrm{II}}(E)\le1$ with
$\Psi_{\mathrm{II}}$ of \eqref{zc9:eq:PsiII}.  Thus if
$\max(\Psi_{\mathrm{deep}},\Psi_{\mathrm{sh}},\Psi_{\mathrm{II}})\le1$ on $E$
then \eqref{zc9:eq:target} holds at every point of $\D$ with $e\in E$ and
$d\le d_{\max}(e)$.

\begin{proposition}[strip]\label{zc9:prop:strip}
For each $m\in\{9,11,13\}$, \eqref{zc9:eq:target} holds at every Region-II point
with $\xi\le1$ and $\tfrac1{60}\le e\le\bar e_m$, where
$\bar e_{13}=\bar e_9=\tfrac{29}{100}$ and $\bar e_{11}=\tfrac{91}{300}$.
\end{proposition}

\begin{proof}
The cover of $[\tfrac1{60},\bar e_m]$ is the \textbf{fixed} rational partition of
Tables~\ref{zc9:tab:m13}--\ref{zc9:tab:m9}: the [ZC] Table-1 breakpoints
$\tfrac1{60},\tfrac1{48},\dots,\tfrac{29}{100}$ (extended by
$\tfrac{3}{10},\tfrac{91}{300}$ for $m=11$), with three of them bisected for
$m=9$ and one for $m=11$ (the bisection midpoints $\tfrac3{160},\tfrac{27}{100},
\tfrac{57}{200}$ are displayed as ordinary rows).  Every row lists an upper
bound (rounded up in the last digit) for the corresponding $\Psi$, an explicit
elementary expression in the rational endpoints; all are $<1$.  On the first few
rows of each table $\Lambda\le\gamma e_L/m$, so the shallow case is empty and
$\Psi_{\mathrm{sh}}$ is marked ``$-$''.  Part~3 recomputes every row; Part~2
independently verifies, on a dense $(e,\kappa)$ grid over each interval, that
the tabulated $\Psi$ of the \emph{applicable} branch dominates the true ratio
$\Def_m/\text{floor}$ (worst slack $+0.05$ at $m=13$).
\end{proof}

\begin{table}[p]\centering\small
\caption{$m=13$ strip: $20$ rows on $[\tfrac1{60},\tfrac{29}{100}]$ (the [ZC] cover, unbisected).}\label{zc9:tab:m13}
\begin{tabular}{llcccl}\toprule
$e_L$&$e_R$&$\Psi_{\mathrm{deep}}\!\le$&$\Psi_{\mathrm{sh}}\!\le$&$\Psi_{\mathrm{II}}\!\le$&bind\\\midrule
$\tfrac{1}{60}$&$\tfrac{1}{48}$&$0.442$&$-$&$0.805$&II\\
$\tfrac{1}{48}$&$\tfrac{1}{40}$&$0.390$&$-$&$0.619$&II\\
$\tfrac{1}{40}$&$\tfrac{1}{32}$&$0.370$&$-$&$0.628$&II\\
$\tfrac{1}{32}$&$\tfrac{1}{25}$&$0.355$&$-$&$0.601$&II\\
$\tfrac{1}{25}$&$\tfrac{1}{20}$&$0.344$&$-$&$0.508$&II\\
$\tfrac{1}{20}$&$\tfrac{1}{16}$&$0.356$&$0.618$&$0.479$&s\\
$\tfrac{1}{16}$&$\tfrac{1}{13}$&$0.337$&$0.477$&$0.447$&s\\
$\tfrac{1}{13}$&$\tfrac{1}{10}$&$0.252$&$0.445$&$0.550$&II\\
$\tfrac{1}{10}$&$\tfrac{1}{8}$&$0.153$&$0.367$&$0.549$&II\\
$\tfrac{1}{8}$&$\tfrac{3}{20}$&$0.095$&$0.335$&$0.435$&II\\
$\tfrac{3}{20}$&$\tfrac{17}{100}$&$0.059$&$0.316$&$0.284$&s\\
$\tfrac{17}{100}$&$\tfrac{9}{50}$&$0.038$&$0.285$&$0.190$&s\\
$\tfrac{9}{50}$&$\tfrac{19}{100}$&$0.034$&$0.315$&$0.164$&s\\
$\tfrac{19}{100}$&$\tfrac{1}{5}$&$0.032$&$0.368$&$0.142$&s\\
$\tfrac{1}{5}$&$\tfrac{21}{100}$&$0.028$&$0.399$&$0.124$&s\\
$\tfrac{21}{100}$&$\tfrac{11}{50}$&$0.022$&$0.378$&$0.108$&s\\
$\tfrac{11}{50}$&$\tfrac{6}{25}$&$0.019$&$0.427$&$0.109$&s\\
$\tfrac{6}{25}$&$\tfrac{13}{50}$&$0.012$&$0.429$&$0.088$&s\\
$\tfrac{13}{50}$&$\tfrac{7}{25}$&$0.009$&$0.481$&$0.077$&s\\
$\tfrac{7}{25}$&$\tfrac{29}{100}$&$0.006$&$0.454$&$0.058$&s\\
\bottomrule\end{tabular}\end{table}
\begin{table}[p]\centering\small
\caption{$m=11$ strip: $23$ rows on $[\tfrac1{60},\tfrac{91}{300}]$ ([ZC] cover $+$ tail $\tfrac3{10}\!\to\!\tfrac{91}{300}$; $[\tfrac1{60},\tfrac1{48}]$ bisected at $\tfrac3{160}$).}\label{zc9:tab:m11}
\begin{tabular}{llcccl}\toprule
$e_L$&$e_R$&$\Psi_{\mathrm{deep}}\!\le$&$\Psi_{\mathrm{sh}}\!\le$&$\Psi_{\mathrm{II}}\!\le$&bind\\\midrule
$\tfrac{1}{60}$&$\tfrac{3}{160}$&$0.496$&$-$&$0.688$&II\\
$\tfrac{3}{160}$&$\tfrac{1}{48}$&$0.462$&$-$&$0.609$&II\\
$\tfrac{1}{48}$&$\tfrac{1}{40}$&$0.450$&$-$&$0.732$&II\\
$\tfrac{1}{40}$&$\tfrac{1}{32}$&$0.425$&$-$&$0.742$&II\\
$\tfrac{1}{32}$&$\tfrac{1}{25}$&$0.403$&$-$&$0.710$&II\\
$\tfrac{1}{25}$&$\tfrac{1}{20}$&$0.385$&$-$&$0.601$&II\\
$\tfrac{1}{20}$&$\tfrac{1}{16}$&$0.393$&$0.807$&$0.566$&s\\
$\tfrac{1}{16}$&$\tfrac{1}{13}$&$0.418$&$0.614$&$0.528$&s\\
$\tfrac{1}{13}$&$\tfrac{1}{10}$&$0.368$&$0.565$&$0.650$&II\\
$\tfrac{1}{10}$&$\tfrac{1}{8}$&$0.242$&$0.457$&$0.650$&II\\
$\tfrac{1}{8}$&$\tfrac{3}{20}$&$0.164$&$0.414$&$0.695$&II\\
$\tfrac{3}{20}$&$\tfrac{17}{100}$&$0.113$&$0.391$&$0.502$&II\\
$\tfrac{17}{100}$&$\tfrac{9}{50}$&$0.079$&$0.354$&$0.365$&II\\
$\tfrac{9}{50}$&$\tfrac{19}{100}$&$0.075$&$0.397$&$0.328$&s\\
$\tfrac{19}{100}$&$\tfrac{1}{5}$&$0.076$&$0.473$&$0.296$&s\\
$\tfrac{1}{5}$&$\tfrac{21}{100}$&$0.070$&$0.520$&$0.268$&s\\
$\tfrac{21}{100}$&$\tfrac{11}{50}$&$0.056$&$0.497$&$0.244$&s\\
$\tfrac{11}{50}$&$\tfrac{6}{25}$&$0.052$&$0.571$&$0.256$&s\\
$\tfrac{6}{25}$&$\tfrac{13}{50}$&$0.037$&$0.585$&$0.227$&s\\
$\tfrac{13}{50}$&$\tfrac{7}{25}$&$0.029$&$0.666$&$0.215$&s\\
$\tfrac{7}{25}$&$\tfrac{29}{100}$&$0.020$&$0.654$&$0.176$&s\\
$\tfrac{29}{100}$&$\tfrac{3}{10}$&$0.020$&$0.787$&$0.189$&s\\
$\tfrac{3}{10}$&$\tfrac{91}{300}$&$0.016$&$0.728$&$0.170$&s\\
\bottomrule\end{tabular}\end{table}
\begin{table}[p]\centering\small
\caption{$m=9$ strip: $23$ rows on $[\tfrac1{60},\tfrac{29}{100}]$ ([ZC] cover; $[\tfrac1{60},\tfrac1{48}],[\tfrac{13}{50},\tfrac7{25}],[\tfrac7{25},\tfrac{29}{100}]$ bisected at $\tfrac3{160},\tfrac{27}{100},\tfrac{57}{200}$).}\label{zc9:tab:m9}
\begin{tabular}{llcccl}\toprule
$e_L$&$e_R$&$\Psi_{\mathrm{deep}}\!\le$&$\Psi_{\mathrm{sh}}\!\le$&$\Psi_{\mathrm{II}}\!\le$&bind\\\midrule
$\tfrac{1}{60}$&$\tfrac{3}{160}$&$0.601$&$-$&$0.841$&II\\
$\tfrac{3}{160}$&$\tfrac{1}{48}$&$0.557$&$-$&$0.745$&II\\
$\tfrac{1}{48}$&$\tfrac{1}{40}$&$0.541$&$-$&$0.894$&II\\
$\tfrac{1}{40}$&$\tfrac{1}{32}$&$0.507$&$-$&$0.907$&II\\
$\tfrac{1}{32}$&$\tfrac{1}{25}$&$0.477$&$-$&$0.867$&II\\
$\tfrac{1}{25}$&$\tfrac{1}{20}$&$0.449$&$-$&$0.734$&II\\
$\tfrac{1}{20}$&$\tfrac{1}{16}$&$0.452$&$-$&$0.693$&II\\
$\tfrac{1}{16}$&$\tfrac{1}{13}$&$0.474$&$0.849$&$0.646$&s\\
$\tfrac{1}{13}$&$\tfrac{1}{10}$&$0.558$&$0.770$&$0.798$&II\\
$\tfrac{1}{10}$&$\tfrac{1}{8}$&$0.396$&$0.613$&$0.803$&II\\
$\tfrac{1}{8}$&$\tfrac{3}{20}$&$0.295$&$0.552$&$0.868$&II\\
$\tfrac{3}{20}$&$\tfrac{17}{100}$&$0.226$&$0.525$&$0.874$&II\\
$\tfrac{17}{100}$&$\tfrac{9}{50}$&$0.173$&$0.479$&$0.691$&II\\
$\tfrac{9}{50}$&$\tfrac{19}{100}$&$0.174$&$0.546$&$0.647$&II\\
$\tfrac{19}{100}$&$\tfrac{1}{5}$&$0.187$&$0.667$&$0.608$&s\\
$\tfrac{1}{5}$&$\tfrac{21}{100}$&$0.183$&$0.744$&$0.574$&s\\
$\tfrac{21}{100}$&$\tfrac{11}{50}$&$0.154$&$0.716$&$0.545$&s\\
$\tfrac{11}{50}$&$\tfrac{6}{25}$&$0.151$&$0.840$&$0.597$&s\\
$\tfrac{6}{25}$&$\tfrac{13}{50}$&$0.118$&$0.874$&$0.574$&s\\
$\tfrac{13}{50}$&$\tfrac{27}{100}$&$0.083$&$0.794$&$0.484$&s\\
$\tfrac{27}{100}$&$\tfrac{7}{25}$&$0.079$&$0.874$&$0.497$&s\\
$\tfrac{7}{25}$&$\tfrac{57}{200}$&$0.068$&$0.869$&$0.467$&s\\
$\tfrac{57}{200}$&$\tfrac{29}{100}$&$0.068$&$0.938$&$0.484$&s\\
\bottomrule\end{tabular}\end{table}


\section{The $m=9$ junction sliver}\label{zc9:sec:sliver}

For $m=9$ the strip stops at $\bar e_9=\tfrac{29}{100}$ (beyond which the
$m$-uniform \emph{shallow} bound $\Psi_{\mathrm{sh}}$ exceeds $1$ near the
Case-I/II boundary) and the corner starts at $e^\star_9=\tfrac{31}{100}$ (its
$\ell$-envelope needs $\delta\le\tfrac7{600}$).  The sliver
$e\in[\tfrac{29}{100},\tfrac{31}{100}]$ (true ratio $\le0.13$ there) is closed by
a \textbf{fixed} $(e,t)$ cell table, $t:=d/\delta$, using the mean-value defect
of Lemma~\ref{zc9:lem:mvt} and the sharper Case-II validity condition below.

\begin{lemma}\label{zc9:lem:certIIval}
The Case-II floor $mK_2fd$ is a valid lower bound for $C_m\psi/(\alpha^3p^{n})$
whenever $2(1+\rho)\xi\ge1$ (weaker than $2\rho\xi>1$).
\end{lemma}
\begin{proof}
The $\lambda=1$ dual (cert-II, [R2]) gives $C_m\psi\ge\sqrt{2\alpha}B_mf
(d-\tfrac{e^2}{16\alpha^2(1+\rho)})$; and
$d-\tfrac{e^2}{16\alpha^2(1+\rho)}\ge\tfrac d2\iff\xi\ge\tfrac1{2(1+\rho)}$; then
relax $\tfrac12\sqrt{2\alpha}B_mf\ge\alpha^3 mK_2f$ as in [ZC]
Lemma~\emph{payfloor}.
\end{proof}

\begin{proposition}[sliver]\label{zc9:prop:sliver}
\eqref{zc9:eq:target} holds for $m=9$ at every Region-II point with
$\tfrac{29}{100}\le e\le\tfrac{31}{100}$ and $d<\delta$.
\end{proposition}

\begin{proof}
First, $R_9\le0$ on $\{t\le\tfrac{16}{25}\}$ throughout
$e\in[\tfrac{29}{100},\tfrac{31}{100}]$ (Part~6: the grid maximum of $R_9$ there
is $-1.4\cdot10^{-6}$); on that region \eqref{zc9:eq:target} is trivial.  On
$t\in[\tfrac{16}{25},1]$ use the fixed $4\times4$ cell table of
Table~\ref{zc9:tab:sliver} (e-strips
$\tfrac{29}{100},\tfrac{59}{200},\tfrac3{10},\tfrac{61}{200},\tfrac{31}{100}$;
$t$-bands $\tfrac{16}{25},\tfrac{73}{100},\tfrac{41}{50},\tfrac{91}{100},1$).
On each cell $\Def_9$ is bounded above by
$\tfrac{n x_U^{\,n-1}}{(\alpha p)^-}\,d_{\mathrm{hi}}(1-\ell_L^{\,m-1})$
(Lemma~\ref{zc9:lem:mvt} with $N\le d(1-\ell^{m-1})$, valid since $2d\le f$ here
by Lemma~\ref{zc9:lem:cornerN}(i), and cell-local monotone envelopes); the
payment is bounded below by $mc_e d_{\mathrm{lo}}^2/e^2$ where $2\rho\xi\le1$ and
by $mK_2fd_{\mathrm{lo}}$ where $2(1+\rho)\xi\ge1$
(Lemma~\ref{zc9:lem:certIIval}), the two conditions covering each cell (their
endpoints of $2(1+\rho)\xi$ are displayed by the script; on any straddling cell
both floors are required, which is what the tabulated ratio uses).  Every
tabulated ratio is $<1$ (worst $0.690$), so $\Def_9\le$ payment on each cell.
\end{proof}

\begin{table}[h]\centering\small
\caption{$m=9$ sliver: $16$ floored $(e,t)$ cells on
$[\tfrac{29}{100},\tfrac{31}{100}]\times[\tfrac{16}{25},1]$; $R_9\le0$ below
$t=\tfrac{16}{25}$.  Ratio $=\Def_9/\text{floor}$ (rounded up).}
\label{zc9:tab:sliver}
\begin{tabular}{lllll}\toprule
$e_L$&$e_R$&$t_L$&$t_R$&ratio $\le$\\\midrule
$\tfrac{29}{100}$&$\tfrac{59}{200}$&$\tfrac{16}{25}$&$\tfrac{73}{100}$&$0.362$\\
$\tfrac{29}{100}$&$\tfrac{59}{200}$&$\tfrac{73}{100}$&$\tfrac{41}{50}$&$0.353$\\
$\tfrac{29}{100}$&$\tfrac{59}{200}$&$\tfrac{41}{50}$&$\tfrac{91}{100}$&$0.346$\\
$\tfrac{29}{100}$&$\tfrac{59}{200}$&$\tfrac{91}{100}$&$1$&$0.340$\\
$\tfrac{59}{200}$&$\tfrac3{10}$&$\tfrac{16}{25}$&$\tfrac{73}{100}$&$0.690$\\
$\tfrac{59}{200}$&$\tfrac3{10}$&$\tfrac{73}{100}$&$\tfrac{41}{50}$&$0.605$\\
$\tfrac{59}{200}$&$\tfrac3{10}$&$\tfrac{41}{50}$&$\tfrac{91}{100}$&$0.343$\\
$\tfrac{59}{200}$&$\tfrac3{10}$&$\tfrac{91}{100}$&$1$&$0.337$\\
$\tfrac3{10}$&$\tfrac{61}{200}$&$\tfrac{16}{25}$&$\tfrac{73}{100}$&$0.000$\\
$\tfrac3{10}$&$\tfrac{61}{200}$&$\tfrac{73}{100}$&$\tfrac{41}{50}$&$0.638$\\
$\tfrac3{10}$&$\tfrac{61}{200}$&$\tfrac{41}{50}$&$\tfrac{91}{100}$&$0.343$\\
$\tfrac3{10}$&$\tfrac{61}{200}$&$\tfrac{91}{100}$&$1$&$0.338$\\
$\tfrac{61}{200}$&$\tfrac{31}{100}$&$\tfrac{16}{25}$&$\tfrac{73}{100}$&$0.000$\\
$\tfrac{61}{200}$&$\tfrac{31}{100}$&$\tfrac{73}{100}$&$\tfrac{41}{50}$&$0.685$\\
$\tfrac{61}{200}$&$\tfrac{31}{100}$&$\tfrac{41}{50}$&$\tfrac{91}{100}$&$0.613$\\
$\tfrac{61}{200}$&$\tfrac{31}{100}$&$\tfrac{91}{100}$&$1$&$0.344$\\
\bottomrule\end{tabular}
\end{table}

\section{The corner $e^\star_m\le e<\tfrac13$}\label{zc9:sec:wide}

Here $\delta=\tfrac{1-3e}6\in(0,\Delta_m]$ with
$\Delta_{13}=\tfrac{13}{600}$, $\Delta_{11}=\tfrac9{600}$,
$\Delta_9=\tfrac7{600}$, $\alpha=\tfrac13+\delta$, $d\in(0,\delta)$ (no use of
$\xi\le1$).  We use the mean-value defect of Lemma~\ref{zc9:lem:mvt} together
with the near-cancellation $d\approx q-L$ forced by $R_m>0$.

\begin{lemma}[corner defect]\label{zc9:lem:cornerN}
For any Region-II point with $0<d<\delta$ (i.e.\ $q>\tfrac13$; in particular on
both the sliver and the corner), with $n=m-2$:
\begin{enumerate}
\item[(i)] $2d\le f$, equivalently $d\le q-L$.
\item[(ii)] If $R_m>0$ then $0<N\le d\,(1-\ell^{\,m-1})$, with $N$ as in
 Lemma~\ref{zc9:lem:mvt}.
\item[(iii)] $\alpha\,\Def_m\le\dfrac{n\,x^{n-1}}p\,N$
 (Lemma~\ref{zc9:lem:mvt}).
\end{enumerate}
\end{lemma}

\begin{proof}
(i) $2d\le f=\alpha-L\iff L\le2q-\alpha$; since $2q-\alpha\ge\tfrac13-\delta>0$,
this is $(2q-\alpha)^2-L^2\ge0$.  With $u:=q-\tfrac13=\delta-d>0$ and
$\alpha=\tfrac13+\delta$, a direct expansion gives
$5q^2-q-4q\alpha+2\alpha^2=u+(5u^2-4u\delta+2\delta^2)>0$
(the quadratic has discriminant $16\delta^2-40\delta^2<0$).  Then $d\le q-L$
because $2d\le\alpha-L=(q-L)+d$.
(ii) $R_m>0\Rightarrow N>0$ (Lemma~\ref{zc9:lem:mvt}); and $d\le q-L$ gives
$N=d-\ell^{m-1}(q-L)\le d-\ell^{m-1}d=d(1-\ell^{m-1})$.
(iii) is Lemma~\ref{zc9:lem:mvt}.
\end{proof}

Let $\alpha(u)=\tfrac{1-u}2$ and, on $[e^\star_m,\tfrac13)$, put
\[
 a_U=\tfrac13+\Delta_m,\quad x_U=\tfrac{a_U}{1-a_U},\quad
 q_L^-=\tfrac23-2\Delta_m,\quad q_L^+=\tfrac23+\tfrac{\Delta_m}2,\quad
 (\alpha p)^-=\tfrac13(1-a_U),
\]
\[
 f_{\!H}=\tfrac{3a_U+\frac13}{q_L^-},\quad E_L=3f_{\!H},\quad
 \ell_L=1-E_L\Delta_m,\quad Q_L=\tfrac{2/3}{q_L^+},\quad
 c^-=\tfrac{2(1-x_U^{m-1})(1-2^{-(m-1)})}{(1+x_U)(3/2)},
\]
\[
 K^-=\tfrac{\sqrt{2/3}\,(1-2^{-(m-1)})}{2a_U^3(3/2)}
\]
(upper bounds for $\alpha,x$; lower/upper bounds for $q+L$; a lower bound for
$\alpha p$; upper bounds for $f/\delta$ and $1-\ell$; a lower bound for $\ell$;
an upper bound for $(q-L)/\delta$; and lower bounds for $c_e,K_2$;
all monotone-endpoint, checked in Part~5).

\begin{proposition}[corner]\label{zc9:prop:wide}
For each $m\in\{9,11,13\}$, \eqref{zc9:eq:target} holds at every Region-II point
with $e^\star_m\le e<\tfrac13$ (indeed all $d<\delta$).
\end{proposition}

\begin{proof}
Assume $R_m>0$.  By Lemma~\ref{zc9:lem:cornerN}(ii)--(iii),
$\Def_m\le\tfrac{n x^{n-1}}{\alpha p}\,d(1-\ell^{m-1})$.
\emph{Case I} ($2\rho\xi\le1$): floor $mc_e\kappa^2\ge9mc_e d^2$
($\kappa=d/e\ge3d$).  Using $d\ge\ell^{m-1}(q-L)\ge\ell_L^{m-1}Q_L\delta$ (from
$R_m>0$ and $q-L\ge Q_L\delta$), $1-\ell^{m-1}\le(m-1)(1-\ell)\le(m-1)E_L\delta$,
$e^2\le\tfrac19$,
\[
 \frac{\Def_m}{mc_e\kappa^2}\ \le\
 \frac{n\,x_U^{\,n-1}}{(\alpha p)^-}\cdot\frac{(m-1)E_L}
 {9\,m\,c^-\,\ell_L^{\,m-1}\,Q_L}\ =:\ r_{\mathrm I}(m).
\]
\emph{Case II} ($2\rho\xi>1$): floor $mK_2fd$; with $f=\alpha(1-\ell)$,
$(1-\ell^{m-1})/f\le(m-1)/\alpha\le3(m-1)$, so
\[
 \frac{\Def_m}{mK_2fd}\ \le\
 \frac{n\,x_U^{\,n-1}}{(\alpha p)^-}\cdot\frac{3(m-1)}{m\,K^-}\ =:\
 r_{\mathrm{II}}(m).
\]
Numerically (Part~5, exact envelopes):
\[
\begin{array}{lccccc}
\toprule
m & \Delta_m & x_U & \ell_L & r_{\mathrm I}(m) & r_{\mathrm{II}}(m)\\
\midrule
13 & \tfrac{13}{600} & 0.5504 & 0.8542 & 0.7063 & 0.0594\\
11 & \tfrac{9}{600}  & 0.5345 & 0.9026 & 0.5893 & 0.1171\\
 9 & \tfrac{7}{600}  & 0.5267 & 0.9256 & 0.9345 & 0.2765\\
\bottomrule
\end{array}
\]
Both are $<1$, so \eqref{zc9:eq:target} holds; if $R_m\le0$ it is trivial.
\end{proof}

\section{Assembly}\label{zc9:sec:assembly}

\begin{proof}[Proof of Theorem~\ref{zc9:thm:main}]
Fix $m\in\{9,11,13\}$ and a Region-II point with $\xi\le1$, $q>\tfrac13$.  If
$e\le\tfrac1{60}$: Corollary~\ref{zc9:cor:za}.  If
$\tfrac1{60}\le e\le\bar e_m$: Proposition~\ref{zc9:prop:strip}.  If
$e\ge e^\star_m$: Proposition~\ref{zc9:prop:wide} (for $m\in\{11,13\}$,
$\bar e_m=e^\star_m$, so the strip meets the corner).  For $m=9$ the sliver
$[\tfrac{29}{100},\tfrac{31}{100}]$ is covered by
Proposition~\ref{zc9:prop:sliver} and $[\tfrac{31}{100},\tfrac13)$ by
Proposition~\ref{zc9:prop:wide}.  These ranges cover $(0,\tfrac13)$.
\end{proof}

\begin{remark}
The independent ground-truth scan (Part~7) gives
$\max R_m/(C_m\psi)=0.122,\,0.094,\,0.080$ for $m=9,11,13$ over the entire
$\xi\le1$ region: the true margins are large everywhere.  The tightest displayed
row is $\Psi_{\mathrm{sh}}=0.938$ (the $m=9$ shallow ridge at $e\approx0.29$) and
$r_{\mathrm I}(9)=0.935$; both reflect slack in the $m$-uniform Case-I
floor/majorant, not in the inequality itself.  The $d_{\mathrm{lo}}$ upgrade of
the previous draft, and its dependence on an inherited bracket, are eliminated:
the only Case-II ingredient beyond [ZC] is the mean-value bound
Lemma~\ref{zc9:lem:mvt}, proved here in full.
\end{remark}
\end{document}
